
\begin{lemma}{7.4.12}\label{lem:7-4-12}
若在归纳的较早阶段构造的
\(
\mathcal M^{\mathrm{main}}_{k',1}
(L;\beta';P'_1,\ldots,P'_{k'})
\)
上的扰动保持循环对称性，
则它们在
\(
\partial\mathcal M^{\mathrm{main}}_{k,1}
(L;\beta;P_1,\ldots,P_k)
\)
上诱导的扰动也保持循环对称性。
\end{lemma}
\begin{proof}
考察如下边界的各分支：
\[
\partial\mathcal M^{\mathrm{main}}_{k,1}
(L;\beta;P_2,\ldots,P_k,P_1)
\]
它们可由下面边界的各分支
\[
\partial\mathcal M^{\mathrm{main}}_{k,1}
(L;\beta;P_1,\ldots,P_k)
\]
中的
\(P_1,\ldots,P_k\)
依次换为
\(P_2,\ldots,P_k,P_1\)
而得到。
对型别~\hyperref[cond:7-4-11-2]{(7.4.11.2)}
中的分支
\[
\mathcal M^{\mathrm{main}}_{k-j+i,1}
\bigl(L;\beta-\beta';
P_1,\ldots,P_{i-1},Q_{i,j}(\beta'),
P_{j+1},\ldots,P_k\bigr)
\]
作此替换，若 \(j<k\)，得到
\begin{equation}
\begin{aligned}
\mathcal M^{\mathrm{main}}_{k-j+i,1}
\bigl(&L;\beta-\beta';
P_2,\ldots,P_i,Q_{i+1,j+1}(\beta'),\\
&P_{j+2},\ldots,P_k,P_1\bigr);
\end{aligned}
\tag{7.4.13.1}\label{eq:7-4-13-1}
\end{equation}
若 \(j=k\)，则得到
\begin{equation}
\mathcal M^{\mathrm{main}}_{i,1}
\bigl(L;\beta-\beta';
P_2,\ldots,P_i,R_{i+1,1}(\beta')\bigr).
\tag{7.4.13.2}\label{eq:7-4-13-2}
\end{equation}

同样，对型别~\hyperref[cond:7-4-11-3]{(7.4.11.3)}
中的分支
\[
\mathcal M^{\mathrm{main}}_{j-i+2,1}
\bigl(L;\beta-\beta';
P_i,\ldots,P_j,R_{j+1,i-1}(\beta')\bigr)
\]
作此替换，若 \(j<k-1\)，得到
\begin{equation}
\begin{aligned}
\mathcal M^{\mathrm{main}}_{j-i+2,1}
\bigl(&L;\beta-\beta';
P_{i+1},\ldots,P_{j+1},\\
&R_{j+2,i}(\beta')\bigr);
\end{aligned}
\tag{7.4.13.3}\label{eq:7-4-13-3}
\end{equation}
若 \(j=k-1\)，则得到
\begin{equation}
\mathcal M^{\mathrm{main}}_{k-i+1,1}
\bigl(L;\beta-\beta';
P_{i+1},\ldots,P_k,Q_{1,i}(\beta')\bigr).
\tag{7.4.13.4}\label{eq:7-4-13-4}
\end{equation}
由归纳假设，
\eqref{eq:7-4-13-1}--\eqref{eq:7-4-13-4}
之和等于
\hyperref[cond:7-4-11-2]{(7.4.11.2)} 和
\hyperref[cond:7-4-11-3]{(7.4.11.3)}
两类分支之和。\translatornote{源书 (7.4.11.3) 允许
\(j=k\)，但这里 (7.4.13.3)
只处理 \(j<k-1\)，
(7.4.13.4) 只处理 \(j=k-1\)，
未说明 \(j=k\) 的变换；
而 (7.4.13.4) 的 \(Q_{1,i}\)
也未直接包含原 \(R_{k,i-1}\)
中 \(P_k\) 的贡献。
此处保留源书读法，待核；
见疑误 A-0559。}
因此后者的总和具有循环对称性。
符号确如上述断言，
可参见第~8.10 节。
型别~\hyperref[cond:7-4-11-1]{(7.4.11.1)}
显然满足结论；
型别~\hyperref[cond:7-4-11-4]{(7.4.11.4)}
则无须另作证明。
\end{proof}

\begin{remark}{7.4.14}\label{rem:7-4-14}
上面引理~\ref{lem:7-4-12}
的证明只用到
\(
\mathcal M^{\mathrm{main}}_{k',1}
(L;\beta';P'_1,\ldots,P'_{k'})
\)
的循环对称性，
不要求
\(
\mathcal M^{\mathrm{main}}_{k',0}
(L;\beta';P'_1,\ldots,P'_{k'})
\)
具有额外的对称性。
\end{remark}

由引理~\ref{lem:7-4-12}，
\(
\mathcal M^{\mathrm{main}}_{k,1}
(L;\beta;P_1,\ldots,P_k)
\)
的边界上已有保持循环对称性的
扰动（多值截面）。
可按如下方式将它延拓到整个模空间，
并保持循环对称性。
先按 \(\mathbb Z_k\) 作用取
\(
\mathcal M^{\mathrm{main}}_{k,1}
(L;\beta;P_1,\ldots,P_k)
\)
的商，得到具有仓西结构的空间
（见定义 A1.45 和引理 A1.49）。
再由定理 A1.23
将商空间边界上的多值截面延拓到内部，
最后将其提升回原模空间，
便得到所需的扰动。

\begin{remark}{7.4.15}\label{rem:7-4-15}
这里说明了为何即使在单调情形下，
也必须使用 \(\mathbb Q\) 系数：
我们需要多值截面在 \(\mathbb Z_k\)
作用下不变，
而这一作用一般有不动点。
\end{remark}

至此，我们已构造出
既保持循环对称性、
又与边界相容的扰动。
由此可按第~\ref{sec:3-8} 节所述
定义算子 \(\mathfrak p\)。
上述边界描述 (7.4.11)
推出
\eqref{eq:3-8-10-1}--\eqref{eq:3-8-10-3}。
为了证明
\eqref{eq:3-8-10-4}--\eqref{eq:3-8-10-6}，
还须确立 \(\mathfrak p\)
与同伦单位的关系。
这只需直接结合上面的论证
与第~\ref{sec:7-3} 节的论证，
故从略。
定理~\ref{thm:3-8-9}
的证明至此完成。

\begin{remark}{7.4.16}\label{rem:7-4-16}
对 \(\mathfrak p\) 实现循环对称性的
上述论证，
比 \(\mathfrak m\) 的情形简单得多。
这是因为
\(\mathfrak p\) 的像是环境流形 \(M\)
上的链，
无须再次对其作用有关的算子
\(\mathfrak p\)。
另一方面，对 \(\mathfrak m\) 而言，
我们必须将其作用两次，
确保所选的 \(L\) 上可数链集满足
\(A_\infty\) 关系
\(\widehat{\mathfrak m}\circ
\widehat{\mathfrak m}=0\)
（见第~\ref{sec:7-2} 节）。
事实上，
第~\ref{sec:7-2} 节为保证
\(A_\infty\) 关系所采用的
\(\mathfrak m\) 构造过程，
势必破坏循环对称性。
\end{remark}
